\documentclass[11pt,twoside]{amsart}
\usepackage[dvipsnames]{xcolor}
\usepackage{hyperref}
\hypersetup{colorlinks=true, linkcolor=Blue, citecolor=Blue}
\usepackage{amsthm, amsmath, amscd, amssymb,centernot,txfonts}
\usepackage{tikz-cd}
\usepackage{lipsum}
\usepackage{cancel} 
\usepackage{multicol}
\usepackage{amsfonts}
\usepackage{amsmath}
\usepackage{mathrsfs}
\usepackage[all]{xy}
\usepackage[left=2.5cm,top=2.5cm,bottom=3cm,right=2.5cm]{geometry}
\usepackage{thmtools}
\usepackage{dirtytalk}
\usepackage{mathtools}
\usetikzlibrary{shapes,arrows}
\usepackage{etoolbox}
\usepackage{dynkin-diagrams}
\usepackage{enumitem}
\usepackage{chngpage}
\usepackage{adjustbox}
\usepackage[normalem]{ulem}
\usepackage[utf8]{inputenc}
\usepackage{fourier} 
\usepackage{array}
\usepackage{makecell}
\usepackage{comment}
%\usepackage[cmtip,all]{xy}
\usepackage{cancel}
\usepackage{cleveref}
\newcommand{\longsquiggly}{\xymatrix{{}\ar@{~>}[r]&{}}}




\setlength{\headheight}{15.2pt}

\renewcommand\theadalign{bc}
\renewcommand\theadfont{\bfseries}
\renewcommand\theadgape{\Gape[4pt]}
\renewcommand\cellgape{\Gape[4pt]}



\setcounter{tocdepth}{1}

\setlength\parskip{0in}
\setlength\parindent{0.2in}


\theoremstyle{plain}
\numberwithin{equation}{section}
\newtheorem{theorem}{Theorem}[section]
\newtheorem{proposition}[theorem]{Proposition}
\newtheorem{lemma}[theorem]{Lemma}
\newtheorem{corollary}[theorem]{Corollary}
\newtheorem{conjecture}[theorem]{Conjecture}
\newtheorem{claim}[theorem]{Claim}
%\newtheorem{set-up}[theorem]{Set-up}

\DeclareMathOperator{\Spec}{Spec}
\DeclareMathOperator{\Proj}{Proj}
%\DeclareMathOperator{\Hom}{Hom}
\DeclareMathOperator{\Der}{Der}
\DeclareMathOperator{\coker}{coker}
%\DeclareMathOperator{\Ext}{\mathcal{E}xt}
\DeclareMathOperator{\im}{im}
\DeclareMathOperator{\rk}{rk}
 
\newcommand{\OO}{\mathcal{O}}
\newcommand{\PP}{\mathbb{P}}
\newcommand{\AAA}{\mathbb{A}}
\newcommand{\CC}{\mathbb{C}}
\newcommand{\ZZ}{\mathbb{Z}}
\newcommand{\Tt}{\mathcal{T}}


\newtheorem{observation}[theorem]{Observation}


\theoremstyle{definition}
\newtheorem{remark}[theorem]{Remark}
\newtheorem{notation}[theorem]{Notation}
\newtheorem{example}[theorem]{Example}
\newtheorem{question}[theorem]{Question}
\newtheorem{set-up}[theorem]{Set-up}
\newtheorem{definition}[theorem]{Definition}
\newtheorem{construction}[definition]{Construction}
\newcommand{\te}{\otimes}

\newcommand{\myeq}{\overset{\mathrm{duality}}{=\joinrel=}}
\newcommand*{\QEDA}{\hfill\ensuremath{\blacksquare}}
\newcommand*{\QEDB}{\hfill\ensuremath{\square}}

\newcommand{\cO}{\mathscr{O}}
\newcommand{\cE}{\mathscr{E}}
\newcommand{\cB}{\mathscr{B}}
\newcommand{\cI}{\mathscr{I}}
\newcommand{\Yt}{\widetilde{Y}}
\newcommand{\Ext}{\operatorname{Ext}}
\newcommand{\Aut}{\operatorname{Aut}}
\newcommand{\Hom}{\operatorname{Hom}}

%\DeclareMathOperator{\Ext}{Ext}
%\DeclareMathOperator{\Hom}{Hom}
%\DeclareMathOperator{\coker}{coker}
\DeclareMathOperator{\Tot}{Tot}
\DeclareMathOperator{\Div}{Div}
\DeclareMathOperator{\Pic}{Pic}
\DeclareMathOperator{\Sing}{Sing}
%\DeclareMathOperator{\Spec}{Spec}
\DeclareMathOperator{\Sym}{Sym}
\DeclareMathOperator{\height}{ht}
 
%\newcommand{\OO}{\mathcal{O}}
\newcommand{\LL}{\mathbf{L}}
\newcommand{\cT}{\mathcal{T}}
\newcommand{\cH}{\mathcal{H}}
%\newcommand{\cI}{\mathcal{I}}
\newcommand{\cJ}{\mathcal{J}}
\newcommand{\cN}{\mathcal{N}}
%\newcommand{\PP}{\mathbb{P}}
%\newcommand{\CC}{\mathbb{C}}
%\newcommand{\ZZ}{\mathbb{Z}}
\newcommand{\HH}{\mathrm{H}}


\tikzstyle{decision} = [diamond, draw, , 
    text width=4.5em, text badly centered, node distance=3cm, inner sep=0pt]
\tikzstyle{block} = [rectangle, draw, , 
    text width=10em, text centered, rounded corners, minimum height=2em]
\tikzstyle{block1} = [rectangle, draw, , 
    text width=5em, text centered, rounded corners, minimum height=2em]    
\tikzstyle{line} = [draw, -latex']
\tikzstyle{cloud} = [draw, ellipse,, node distance=3cm,
    minimum height=2em]

\begin{document}

\title[Deformation theory of morphisms finite onto a singular base ]{Deformation theory of morphisms finite onto a singular base}

\author[J. Mukherjee]{Jayan Mukherjee}
\address{Department of Mathematics, Oklahoma State University, Stillwater, USA}
\email{jayan.mukherjee@okstate.edu}

\maketitle


\section{Reflexive Ribbons}

\begin{definition}[Ropes with conormal sheaf]
\label{def:ribbon-conormal-sheaf}
Let \(Y\) be a reduced connected scheme, and let \(\mathcal E\) be a coherent
\(\mathcal O_Y\)-module. A \emph{ribbon over \(Y\)
with conormal sheaf \(\mathcal E\)} is a scheme \(\widetilde Y\), together with a closed
immersion
\[
    Y\hookrightarrow \widetilde Y,
\]
such that
\[
    \widetilde Y_{\mathrm{red}}=Y,
    \qquad
    \mathcal I_{Y/\widetilde Y}^{2}=0,
    \qquad
    \mathcal I_{Y/\widetilde Y}\simeq \mathcal E
\]
as \(\mathcal O_Y\)-modules. Equivalently, \(\widetilde Y\) fits into an exact sequence of sheaves of rings
\[
    0\longrightarrow \mathcal E
    \longrightarrow \mathcal O_{\widetilde Y}
    \longrightarrow \mathcal O_Y
    \longrightarrow 0,
\]
where \(\mathcal E\) is an ideal of square zero. If \(\mathcal E\) is reflexive of rank \(n-1\), then $\widetilde{Y}$ is called a \emph{reflexive \(n\)-rope} over \(Y\).
\end{definition}



\begin{remark}
\label{rem:properties-reflexive-rope}
Assume that \(Y\) is normal and that \(\mathcal E\) is reflexive.

\begin{enumerate}
    \item Since \(Y\) is normal, a reflexive sheaf on \(Y\) is locally free outside a closed
    subset of codimension at least two. Hence a reflexive rope is an ordinary rope on
    the open subset where \(\mathcal E\) is locally free. In particular, a reflexive rope
    is a genuine rope in codimension one.

    \item The scheme \(\widetilde Y\) is a square-zero thickening of \(Y\). More precisely,
    the closed immersion \(Y\hookrightarrow \widetilde Y\) is defined by the ideal
    \[
        \mathcal I_{Y/\widetilde Y}\simeq \mathcal E,
    \]
    and this ideal satisfies
    \[
        \mathcal I_{Y/\widetilde Y}^{2}=0.
    \]
    Thus the underlying topological space of \(\widetilde Y\) is the same as that of \(Y\),
    but the structure sheaf has an additional square-zero layer given by \(\mathcal E\).

    \item If \(\mathcal E\) has rank \(n-1\), then \(\widetilde Y\) has generic multiplicity
    \(n\) along \(Y\). Indeed, at the generic point \(\eta\) of \(Y\), one has
    \[
        \operatorname{length}_{\mathcal O_{Y,\eta}}
        \mathcal O_{\widetilde Y,\eta}
        =
        1+\operatorname{rank}\mathcal E
        =
        n.
    \]
    In particular, when \(\mathcal E\) has rank one, \(\widetilde Y\) is generically a
    double structure over \(Y\).

    \item \textcolor{blue}{Since \(Y\) is normal, \(\mathcal O_Y\) satisfies Serre's condition \(S_2\).
    Since \(\mathcal E\) is reflexive, \(\mathcal E\) also satisfies \(S_2\). From the exact
    sequence
    \[
        0\longrightarrow \mathcal E
        \longrightarrow \mathcal O_{\widetilde Y}
        \longrightarrow \mathcal O_Y
        \longrightarrow 0,
    \]
    the depth lemma implies that \(\mathcal O_{\widetilde Y}\) satisfies \(S_2\). Hence
    \(\widetilde Y\) is an \(S_2\)-scheme.}

    \item \textcolor{blue}{Reflexivity alone does not imply that \(\widetilde Y\) is Cohen--Macaulay in
    dimensions at least three. If \(Y\) is Cohen--Macaulay and \(\mathcal E\) is maximal
    Cohen--Macaulay as an \(\mathcal O_Y\)-module, then the same exact sequence shows
    that \(\widetilde Y\) is Cohen--Macaulay.}

   
\end{enumerate}
\end{remark}



\begin{proposition}[Injectivity of the conormal map for torsion-free square-zero ideals]
\label{prop:injectivity-conormal-map}
Let \(k\) be a field of characteristic \(0\). Let \(A\) be a reduced noetherian
\(k\)-algebra, and let
\[
    0\longrightarrow M\longrightarrow B\xrightarrow{q} A\longrightarrow 0
\]
be a square-zero extension of \(k\)-algebras; that is,
\[
    M=\ker(q)
    \qquad\text{and}\qquad
    M^2=0.
\]
Assume that \(M\) is torsion-free as an \(A\)-module, where torsion-free means that
multiplication by every nonzerodivisor of \(A\) is injective on \(M\). Then the natural
conormal map
\[
    M=M/M^2
    \longrightarrow
    \Omega_{B/k}\otimes_B A
\]
is injective.

Equivalently, we have an exact sequence the conormal sequence
\[
  0 \longrightarrow  M
    \longrightarrow
    \Omega_{B/k}\otimes_B A
    \longrightarrow
    \Omega_{A/k}
    \longrightarrow 0
\]
is left-exact.
\end{proposition}

\begin{proof}
Let
\[
    \delta:M\longrightarrow \Omega_{B/k}\otimes_B A
\]
be the conormal map. Explicitly, it is given by
\[
    m\longmapsto dm\otimes 1.
\]
Let
\[
    N:=\ker(\delta).
\]
We want to prove that \(N=0\).

We first reduce to the generic points of \(\operatorname{Spec}A\). Since \(A\) is noetherian,
it has only finitely many minimal primes; see \cite[Algebra, Lemma 10.31.6]{stacks-project}.
Let \(\mathfrak p\subset A\) be a minimal prime. Since \(A\) is reduced, the localization
\[
    A_{\mathfrak p}
\]
is a field; see \cite[Algebra, Lemma 10.25.1]{stacks-project}. We denote this field by
\[
    K:=A_{\mathfrak p}.
\]

Localizing the given square-zero extension at \(\mathfrak p\), we obtain an exact sequence
\[
    0\longrightarrow M_{\mathfrak p}
    \longrightarrow B_{\mathfrak p}
    \longrightarrow K
    \longrightarrow 0.
\]
The ideal \(M_{\mathfrak p}\) is still square-zero.

Because \(\operatorname{char}k=0\), every field extension of \(k\) is formally smooth over
\(k\); see \cite[Algebra, Lemma 10.158.8]{stacks-project}. Hence \(K/k\) is formally
smooth. Applying the infinitesimal lifting property for formal smoothness to the square-zero
surjection
\[
    B_{\mathfrak p}\twoheadrightarrow K
\]
and to the identity map \(K\to K\), we obtain a \(k\)-algebra section
\[
    K\longrightarrow B_{\mathfrak p}.
\]

Now apply the split conormal exact sequence for Kähler differentials to the surjection
\[
    B_{\mathfrak p}\twoheadrightarrow K
\]
with kernel \(M_{\mathfrak p}\). The split conormal sequence says that, since this surjection
has a \(k\)-algebra section, the sequence
\[
    0
    \longrightarrow
    M_{\mathfrak p}/M_{\mathfrak p}^{2}
    \longrightarrow
    \Omega_{B_{\mathfrak p}/k}\otimes_{B_{\mathfrak p}} K
    \longrightarrow
    \Omega_{K/k}
    \longrightarrow 0
\]
is exact; see \cite[Algebra, Lemma 10.131.10]{stacks-project}. Since
\[
    M_{\mathfrak p}^{2}=0,
\]
we have
\[
    M_{\mathfrak p}/M_{\mathfrak p}^{2}=M_{\mathfrak p}.
\]
Therefore the localized conormal map
\[
    M_{\mathfrak p}
    \longrightarrow
    \Omega_{B_{\mathfrak p}/k}\otimes_{B_{\mathfrak p}} K
\]
is injective.

Formation of Kähler differentials commutes with localization; see
\cite[Algebra, Lemma 10.131.8]{stacks-project}. Hence
\[
    \Omega_{B_{\mathfrak p}/k}
    \simeq
    \Omega_{B/k}\otimes_B B_{\mathfrak p}.
\]
After tensoring with \(K=A_{\mathfrak p}\), this identifies
\[
    \Omega_{B_{\mathfrak p}/k}\otimes_{B_{\mathfrak p}} K
\]
with
\[
    \left(\Omega_{B/k}\otimes_B A\right)_{\mathfrak p}.
\]
Thus the localized map
\[
    \delta_{\mathfrak p}:M_{\mathfrak p}
    \longrightarrow
    \left(\Omega_{B/k}\otimes_B A\right)_{\mathfrak p}
\]
is injective. Equivalently,
\[
    N_{\mathfrak p}=0
\]
for every minimal prime \(\mathfrak p\subset A\).

It remains to pass from generic injectivity to global injectivity. Let \(n\in N\). Since
\(N_{\mathfrak p}=0\) for every minimal prime \(\mathfrak p\), for each minimal prime
\(\mathfrak p\) there exists an element
\[
    s_{\mathfrak p}\in A\setminus \mathfrak p
\]
such that
\[
    s_{\mathfrak p}n=0.
\]
Equivalently, the annihilator ideal
\[
    \operatorname{Ann}_A(n)
\]
is not contained in any minimal prime of \(A\).

Since \(A\) has only finitely many minimal primes, prime avoidance
\cite[Algebra, Lemma 10.15.2]{stacks-project} implies that
\[
    \operatorname{Ann}_A(n)
\]
is not contained in the union of the minimal primes. Hence there exists an element
\[
    s\in \operatorname{Ann}_A(n)
\]
which is contained in no minimal prime of \(A\).

Because \(A\) is reduced, the set of zerodivisors of \(A\) is the union of its minimal
primes; see \cite[Algebra, Lemma 10.25.2]{stacks-project}. Therefore \(s\) is a
nonzerodivisor of \(A\). Since \(M\) is torsion-free and \(N\subset M\), multiplication by
\(s\) is injective on \(N\). But
\[
    sn=0
\]
by construction. Hence
\[
    n=0.
\]
Therefore every element of \(N\) is zero, so
\[
    N=0.
\]

Thus the conormal map
\[
    M\longrightarrow \Omega_{B/k}\otimes_B A
\]
is injective.
\end{proof}


\begin{proposition}[Classification of square-zero thickenings with torsion-free conormal sheaf]
\label{prop:classification-square-zero-thickenings}
Let \(k\) be a field of characteristic \(0\). Let \(Y\) be a reduced connected
noetherian \(k\)-scheme, and let \(\mathcal E\) be a coherent torsion-free
\(\mathcal O_Y\)-module. Consider pairs
\[
    (j:Y\hookrightarrow \widetilde Y,\alpha),
\]
where \(j\) is a closed immersion such that
\[
    \widetilde Y_{\mathrm{red}}=Y,
    \qquad
    \mathcal I_{Y/\widetilde Y}^{2}=0,
\]
and where
\[
    \alpha:\mathcal I_{Y/\widetilde Y}\xrightarrow{\sim}\mathcal E
\]
is an isomorphism of \(\mathcal O_Y\)-modules. Then the isomorphism classes of such pairs, where isomorphisms are required to be
over \(Y\) and compatible with the chosen identification with \(\mathcal E\), are
classified by
\[
    \operatorname{Ext}^{1}_{Y}(\Omega_{Y/k},\mathcal E).
\]

If the identification
\[
    \mathcal I_{Y/\widetilde Y}\simeq\mathcal E
\]
is not fixed, then the set of isomorphism classes is the quotient of
\[
    \operatorname{Ext}^{1}_{Y}(\Omega_{Y/k},\mathcal E)
\]
by the natural action of
\[
    \operatorname{Aut}_{\mathcal O_Y}(\mathcal E).
\]
\end{proposition}

\begin{proof}
Let
\[
    (j:Y\hookrightarrow \widetilde Y,\alpha)
\]
be such a square-zero thickening, and set
\[
    \mathcal K:=\mathcal I_{Y/\widetilde Y}.
\]
Thus
\[
    \mathcal K^{2}=0
\]
and
\[
    \alpha:\mathcal K\xrightarrow{\sim}\mathcal E.
\]

The closed immersion
\[
    j:Y\hookrightarrow \widetilde Y
\]
has the usual conormal sequence
\[
    \mathcal K/\mathcal K^{2}
    \longrightarrow
    \Omega_{\widetilde Y/k}\otimes_{\mathcal O_{\widetilde Y}}\mathcal O_Y
    \longrightarrow
    \Omega_{Y/k}
    \longrightarrow 0.
\]
Since \(\mathcal K^2=0\), this becomes
\[
    \mathcal K
    \longrightarrow
    \Omega_{\widetilde Y/k}|_Y
    \longrightarrow
    \Omega_{Y/k}
    \longrightarrow 0.
\]

We claim that the first map is injective. This is local on \(Y\). On an affine open
subset \(U=\operatorname{Spec}A\subset Y\), write
\[
    \widetilde U=\operatorname{Spec}B
\]
for the corresponding affine open subset of \(\widetilde Y\). Then we have a
square-zero extension
\[
    0\longrightarrow M\longrightarrow B\longrightarrow A\longrightarrow 0,
\]
where
\[
    M=\Gamma(U,\mathcal K).
\]
Since \(\mathcal K\simeq\mathcal E\), and since \(\mathcal E\) is torsion-free, the
\(A\)-module \(M\) is torsion-free. Also \(A\) is reduced because \(Y\) is reduced.
Therefore, by \Cref{prop:injectivity-conormal-map}, the conormal map
\[
    M
    \longrightarrow
    \Omega_{B/k}\otimes_B A
\]
is injective. Hence the sheaf map
\[
    \mathcal K
    \longrightarrow
    \Omega_{\widetilde Y/k}|_Y
\]
is injective.

Thus we obtain a short exact sequence
\[
    0
    \longrightarrow
    \mathcal K
    \longrightarrow
    \Omega_{\widetilde Y/k}|_Y
    \longrightarrow
    \Omega_{Y/k}
    \longrightarrow 0.
\]
Using the chosen isomorphism
\[
    \alpha:\mathcal K\xrightarrow{\sim}\mathcal E,
\]
this gives a short exact sequence
\[
    0
    \longrightarrow
    \mathcal E
    \longrightarrow
    \Omega_{\widetilde Y/k}|_Y
    \longrightarrow
    \Omega_{Y/k}
    \longrightarrow 0.
\]
Therefore the thickening determines an extension class
\[
    [e_{\widetilde Y}]
    \in
    \operatorname{Ext}^{1}_{Y}(\Omega_{Y/k},\mathcal E).
\]

Conversely, let
\[
    0
    \longrightarrow
    \mathcal E
    \longrightarrow
    \mathcal G
    \xrightarrow{p}
    \Omega_{Y/k}
    \longrightarrow 0
\]
be an extension of \(\Omega_{Y/k}\) by \(\mathcal E\). We construct from this a
square-zero thickening of \(Y\) with conormal sheaf \(\mathcal E\).

Let
\[
    d:\mathcal O_Y\longrightarrow \Omega_{Y/k}
\]
be the universal derivation. Define
\[
    \mathcal O_{\widetilde Y}
    :=
    \mathcal O_Y\times_{\Omega_{Y/k}}\mathcal G.
\]
Thus, for an open set \(U\subset Y\),
\[
    \mathcal O_{\widetilde Y}(U)
    =
    \left\{
        (f,g)\in \mathcal O_Y(U)\oplus\mathcal G(U)
        \;\middle|\;
        df=p(g)
    \right\}.
\]

We define multiplication on \(\mathcal O_{\widetilde Y}\) by
\[
    (f,g)\cdot(f',g')
    :=
    (ff',fg'+f'g).
\]
This multiplication is well-defined. Indeed, since
\[
    p(g)=df
    \qquad\text{and}\qquad
    p(g')=df',
\]
we have
\[
    p(fg'+f'g)
    =
    f\,p(g')+f'\,p(g)
    =
    f\,df'+f'\,df
    =
    d(ff').
\]
Thus
\[
    (ff',fg'+f'g)
\]
again lies in the fiber product
\[
    \mathcal O_Y\times_{\Omega_{Y/k}}\mathcal G.
\]
With this multiplication, \(\mathcal O_{\widetilde Y}\) is a sheaf of rings whose unit is
\[
    (1,0).
\]

The projection onto the first factor gives a surjection
\[
    \mathcal O_{\widetilde Y}
    \longrightarrow
    \mathcal O_Y,
    \qquad
    (f,g)\longmapsto f.
\]
Its kernel consists of pairs
\[
    (0,g)
\]
with
\[
    p(g)=0.
\]
Since
\[
    \ker(p)=\mathcal E,
\]
we get
\[
    \ker(\mathcal O_{\widetilde Y}\to\mathcal O_Y)\simeq\mathcal E.
\]
Moreover, if \(e,e'\) are local sections of \(\mathcal E\), then
\[
    (0,e)\cdot(0,e')=(0,0).
\]
Hence this kernel is an ideal of square zero.

Therefore the ringed space defined by \(\mathcal O_{\widetilde Y}\) is a square-zero
thickening
\[
    Y\hookrightarrow \widetilde Y
\]
whose ideal sheaf is identified with \(\mathcal E\).

These two constructions give the desired classification of marked square-zero thickenings
by
\[
    \operatorname{Ext}^{1}_{Y}(\Omega_{Y/k},\mathcal E).
\]

Finally, suppose the identification of the square-zero ideal with \(\mathcal E\) is not fixed.
Changing an identification
\[
    \mathcal I_{Y/\widetilde Y}\xrightarrow{\sim}\mathcal E
\]
by an automorphism
\[
    \sigma\in \operatorname{Aut}_{\mathcal O_Y}(\mathcal E)
\]
pushes out the corresponding extension class by \(\sigma\). Thus
\[
    \operatorname{Aut}_{\mathcal O_Y}(\mathcal E)
\]
acts naturally on
\[
    \operatorname{Ext}^{1}_{Y}(\Omega_{Y/k},\mathcal E),
\]
and the unmarked isomorphism classes are obtained by taking the quotient by this action.
\end{proof}



\begin{proposition}[Morphisms from square-zero thickenings to an ambient scheme]
\label{prop:square-zero-thickenings-and-morphisms}
Let
\[
    i:Y\hookrightarrow Z
\]
be a closed immersion of schemes, and let
\[
    I=I_{Y/Z}\subset \mathcal O_Z
\]
be its ideal sheaf. Let \(\mathcal E\) be an \(\mathcal O_Y\)-module. Consider triples
\[
    (\widetilde Y,j,\alpha),
\]
where
\[
    j:Y\hookrightarrow \widetilde Y
\]
is a closed immersion defined by a square-zero ideal
\[
    \mathcal K:=\mathcal I_{Y/\widetilde Y},
    \qquad
    \mathcal K^2=0,
\]
and
\[
    \alpha:\mathcal K\xrightarrow{\sim}\mathcal E
\]
is an isomorphism of \(\mathcal O_Y\)-modules.

Then the following hold.

\begin{enumerate}
    \item Isomorphism classes of pairs
    \[
        \left((\widetilde Y,j,\alpha),\widetilde i\right),
    \]
    where
    \[
        \widetilde i:\widetilde Y\to Z
    \]
    is a morphism satisfying
    \[
        \widetilde i\circ j=i,
    \]
    are naturally in bijection with homomorphisms
    \[
        \tau:I/I^2\longrightarrow \mathcal E.
    \]

    \item Under this correspondence, the scheme-theoretic image
    \[
        W:=\operatorname{Im}(\widetilde i)\subset Z
    \]
    has ideal sheaf
    \[
        I_{W/Z}
        =
        \ker\left(
            I
            \longrightarrow
            I/I^2
            \xrightarrow{\tau}
            \mathcal E
        \right).
    \]
    Consequently,
    \[
        I^2\subset I_{W/Z}\subset I,
    \]
    and the ideal of \(Y\) inside \(W\) is
    \[
        \mathcal I_{Y/W}
        \simeq
        \operatorname{im}(\tau).
    \]

    \item The morphism
    \[
        \widetilde i:\widetilde Y\to Z
    \]
    is a closed immersion if and only if
    \[
        \tau:I/I^2\to \mathcal E
    \]
    is surjective.

    In particular, embedded square-zero thickenings
    \[
        \widetilde Y\subset Z
    \]
    with
    \[
        \mathcal I_{Y/\widetilde Y}\simeq\mathcal E
    \]
    correspond precisely to surjections
    \[
        I/I^2\twoheadrightarrow \mathcal E.
    \]
\end{enumerate}
\end{proposition}

\begin{proof}
We first explain how a morphism
\[
    \widetilde i:\widetilde Y\to Z
\]
extending \(i:Y\hookrightarrow Z\) produces a homomorphism
\[
    \tau:I/I^2\to \mathcal E.
\]

Let
\[
    \mathcal K=\mathcal I_{Y/\widetilde Y}.
\]
Since \(\widetilde i\circ j=i\), the morphism on structure sheaves
\[
    \widetilde i^\#:\mathcal O_Z\longrightarrow \widetilde i_*\mathcal O_{\widetilde Y}
\]
reduces modulo \(\mathcal K\) to the map
\[
    i^\#:\mathcal O_Z\longrightarrow i_*\mathcal O_Y.
\]
Therefore, if \(a\) is a local section of \(I\), then \(i^\#(a)=0\), and hence
\[
    \widetilde i^\#(a)\in \mathcal K.
\]
Using the fixed isomorphism
\[
    \alpha:\mathcal K\xrightarrow{\sim}\mathcal E,
\]
we define
\[
    \lambda_{\widetilde i}(a)
    :=
    \alpha\bigl(\widetilde i^\#(a)\bigr)
    \in \mathcal E.
\]

We claim that \(\lambda_{\widetilde i}\) kills \(I^2\). Indeed, if \(a,b\) are local
sections of \(I\), then
\[
    \widetilde i^\#(a),\widetilde i^\#(b)\in \mathcal K.
\]
Since \(\mathcal K^2=0\), it follows that
\[
    \widetilde i^\#(ab)
    =
    \widetilde i^\#(a)\widetilde i^\#(b)
    =
    0.
\]
Thus
\[
    \lambda_{\widetilde i}(ab)=0.
\]
Therefore \(\lambda_{\widetilde i}\) descends to a homomorphism
\[
    \tau_{\widetilde i}:I/I^2\longrightarrow \mathcal E.
\]

This homomorphism is \(\mathcal O_Y\)-linear. To see this, let \(f\) be a local section
of \(\mathcal O_Z\), and let \(a\) be a local section of \(I\). Then
\[
    \widetilde i^\#(fa)
    =
    \widetilde i^\#(f)\widetilde i^\#(a).
\]
The section \(\widetilde i^\#(a)\) lies in \(\mathcal K\). Since \(\mathcal K^2=0\), the
\(\mathcal O_{\widetilde Y}\)-module structure on \(\mathcal K\) factors through
\[
    \mathcal O_{\widetilde Y}/\mathcal K\simeq\mathcal O_Y.
\]
Thus
\[
    \widetilde i^\#(f)\widetilde i^\#(a)
    =
    \overline f\cdot \widetilde i^\#(a),
\]
where \(\overline f\) is the image of \(f\) in \(\mathcal O_Y\). Applying \(\alpha\), we get
\[
    \tau_{\widetilde i}(\overline{fa})
    =
    \overline f\cdot \tau_{\widetilde i}(\overline a).
\]
Hence
\[
    \tau_{\widetilde i}:I/I^2\to\mathcal E
\]
is \(\mathcal O_Y\)-linear.

We now construct the inverse correspondence.

Let
\[
    \tau:I/I^2\to\mathcal E
\]
be an \(\mathcal O_Y\)-module homomorphism. Let
\[
    Y^{(1)}_Z:=V(I^2)\subset Z
\]
be the first infinitesimal neighborhood of \(Y\) in \(Z\). Its structure sheaf fits into
the square-zero extension
\[
    0\longrightarrow I/I^2
    \longrightarrow \mathcal O_{Y^{(1)}_Z}
    \longrightarrow \mathcal O_Y
    \longrightarrow 0.
\]
We push out this square-zero extension along
\[
    \tau:I/I^2\to\mathcal E.
\]

Concretely, define a sheaf of rings \(\mathcal O_{\widetilde Y_\tau}\) on the underlying
topological space of \(Y\) by
\[
    \mathcal O_{\widetilde Y_\tau}
    :=
    \left(
        \mathcal O_{Y^{(1)}_Z}\oplus \mathcal E
    \right)
    \big/
    \left\langle
        (u,-\tau(u))\mid u\in I/I^2
    \right\rangle.
\]
Here \(I/I^2\) is viewed as a square-zero ideal in \(\mathcal O_{Y^{(1)}_Z}\), and
\(\mathcal E\) is placed as a square-zero ideal. The multiplication is the one induced
from the standard square-zero multiplication:
\[
    (a,e)(a',e')
    =
    \bigl(
        aa',
        \overline a\, e'
        +
        \overline{a'}\, e
    \bigr),
\]
where \(\overline a,\overline{a'}\) denote the images in \(\mathcal O_Y\).

By construction, there is an exact sequence of sheaves of rings and modules
\[
    0\longrightarrow \mathcal E
    \longrightarrow \mathcal O_{\widetilde Y_\tau}
    \longrightarrow \mathcal O_Y
    \longrightarrow 0,
\]
and \(\mathcal E\) is an ideal of square zero. Therefore
\[
    j_\tau:Y\hookrightarrow \widetilde Y_\tau
\]
is a square-zero thickening with conormal sheaf \(\mathcal E\).

The natural quotient map
\[
    \mathcal O_Z
    \longrightarrow
    \mathcal O_Z/I^2
    =
    \mathcal O_{Y^{(1)}_Z}
\]
followed by
\[
    \mathcal O_{Y^{(1)}_Z}\to \mathcal O_{\widetilde Y_\tau}
\]
defines a morphism
\[
    \widetilde i_\tau:\widetilde Y_\tau\to Z.
\]
Modulo the square-zero ideal \(\mathcal E\), this morphism restricts to the original
closed immersion
\[
    i:Y\hookrightarrow Z.
\]
Hence
\[
    \widetilde i_\tau\circ j_\tau=i.
\]

We now check that the two constructions are inverse to each other.

Starting from a homomorphism \(\tau:I/I^2\to\mathcal E\), the construction above gives
\[
    \widetilde i_\tau:\widetilde Y_\tau\to Z.
\]
If \(a\) is a local section of \(I\), then its image in
\(\mathcal O_{Y^{(1)}_Z}\) is the class \(\overline a\in I/I^2\). In the pushout
\(\mathcal O_{\widetilde Y_\tau}\), this class becomes \(\tau(\overline a)\in\mathcal E\).
Therefore the homomorphism associated to \(\widetilde i_\tau\) is exactly \(\tau\).

Conversely, suppose we start with a pair
\[
    \left((\widetilde Y,j,\alpha),\widetilde i\right).
\]
The morphism
\[
    \widetilde i^\#:\mathcal O_Z\to\widetilde i_*\mathcal O_{\widetilde Y}
\]
kills \(I^2\). Hence it
factors through
\[
    \mathcal O_{Y^{(1)}_Z}=\mathcal O_Z/I^2.
\]
We obtain a morphism of square-zero extensions
\[
\begin{tikzcd}
    0 \arrow[r] &
    I/I^2 \arrow[r] \arrow[d,"\tau_{\widetilde i}"] &
    \mathcal O_{Y^{(1)}_Z} \arrow[r] \arrow[d] &
    \mathcal O_Y \arrow[r] \arrow[d,equal] &
    0
    \\
    0 \arrow[r] &
    \mathcal E \arrow[r] &
    \mathcal O_{\widetilde Y} \arrow[r] &
    \mathcal O_Y \arrow[r] &
    0 .
\end{tikzcd}
\]
By the universal property of the pushout, this induces a morphism
\[
    \mathcal O_{\widetilde Y_{\tau_{\widetilde i}}}
    \longrightarrow
    \mathcal O_{\widetilde Y}.
\]
This morphism is the identity on the quotient \(\mathcal O_Y\) and the identity on the
square-zero ideal \(\mathcal E\). Hence it is an isomorphism of sheaves of rings. Therefore
the original pair is recovered from \(\tau_{\widetilde i}\). This proves the bijection in
part \((1)\).

Next we compute the ideal of the scheme-theoretic image. Let
\[
    W=\operatorname{Im}(\widetilde i_\tau)\subset Z.
\]
By definition,
\[
    I_{W/Z}
    =
    \ker\left(
        \mathcal O_Z\to \mathcal O_{\widetilde Y_\tau}
    \right).
\]
If a local section \(f\in\mathcal O_Z\) maps to zero in \(\mathcal O_{\widetilde Y_\tau}\),
then its image in \(\mathcal O_Y\) is zero, so \(f\in I\). For \(a\in I\), its image in
\(\mathcal O_{\widetilde Y_\tau}\) is precisely
\[
    \tau(\overline a)\in\mathcal E,
\]
where \(\overline a\) denotes the class of \(a\) in \(I/I^2\). Therefore
\[
    a\in I_{W/Z}
    \quad\Longleftrightarrow\quad
    \tau(\overline a)=0.
\]
Thus
\[
    I_{W/Z}
    =
    \ker\left(
        I
        \longrightarrow
        I/I^2
        \xrightarrow{\tau}
        \mathcal E
    \right).
\]

Since this ideal contains \(I^2\) and is contained in \(I\), we have
\[
    I^2\subset I_{W/Z}\subset I.
\]
Moreover, the ideal of \(Y\) in \(W\) is
\[
    \mathcal I_{Y/W}=I/I_{W/Z}.
\]
By the first isomorphism theorem,
\[
    I/I_{W/Z}\simeq\operatorname{im}(\tau).
\]
Hence
\[
    \mathcal I_{Y/W}\simeq \operatorname{im}(\tau).
\]

Finally, we prove the closed immersion criterion. The morphism
\[
    \widetilde i_\tau:\widetilde Y_\tau\to Z
\]
is a closed immersion if and only if the corresponding morphism of sheaves of rings
\[
    \mathcal O_Z\to \mathcal O_{\widetilde Y_\tau}
\]
is surjective. Since
\[
    \mathcal O_Z\to\mathcal O_Y
\]
is already surjective, the only possible failure of surjectivity lies in the square-zero
ideal \(\mathcal E\subset\mathcal O_{\widetilde Y_\tau}\). The image of \(I\) inside
\(\mathcal E\) is exactly
\[
    \operatorname{im}(\tau).
\]
Therefore
\[
    \mathcal O_Z\to \mathcal O_{\widetilde Y_\tau}
\]
is surjective if and only if
\[
    \operatorname{im}(\tau)=\mathcal E,
\]
i.e. if and only if
\[
    \tau:I/I^2\to\mathcal E
\]
is surjective.

Thus \(\widetilde i_\tau\) is a closed immersion if and only if \(\tau\) is surjective.
This proves the proposition.
\end{proof}


\section{The map connecting deformations and reflexive ribbons}

\begin{lemma}[The induced map to the two conormal Hom groups]
\label{lem:T1-map-to-conormal-Hom}
Let \(k\) be an algebraically closed field, and let
\[
    X \xrightarrow{\pi} Y \xrightarrow{i} Z
\]
be morphisms of \(k\)-schemes, with \(\pi\) finite and \(i\) a closed immersion.
Let \(I=I_{Y/Z}\) be the ideal sheaf of \(Y\) in \(Z\), and set
\[
    \varphi=i\circ \pi.
\]
Assume that the finite \(\mathcal O_Y\)-algebra \(\pi_*\mathcal O_X\) admits a splitting
\[
    \pi_*\mathcal O_X \simeq \mathcal O_Y\oplus \mathcal E
\]
as an \(\mathcal O_Y\)-module. Then there is a natural homomorphism
\[
    \Theta=\Theta_1\oplus \Theta_2:
    T^1(X/Z)
    \longrightarrow
    \operatorname{Hom}_{\mathcal O_Y}(I/I^2,\mathcal O_Y)
    \oplus
    \operatorname{Hom}_{\mathcal O_Y}(I/I^2,\mathcal E),
\]
where
\[
    T^1(X/Z):=\operatorname{Ext}^1_X(L_{X/Z},\mathcal O_X).
\]
Equivalently, \(\Theta\) is the composite
\[
\begin{tikzcd}[column sep=large]
    T^1(X/Z)
    \arrow[r]
    &
    \operatorname{Ext}^1_X(L\pi^*L_{Y/Z},\mathcal O_X)
    \arrow[r,"\sim"]
    &
    \operatorname{Hom}_{\mathcal O_Y}(I/I^2,\pi_*\mathcal O_X)
    \arrow[r,"\sim"]
    &
    \operatorname{Hom}(I/I^2,\mathcal O_Y)
    \oplus
    \operatorname{Hom}(I/I^2,\mathcal E).
\end{tikzcd}
\]
\end{lemma}

\begin{proof}
We use the cotangent complex. By Wehler's convention for tangent cohomology,
\[
    T^i(X/Z,\mathcal M)
    =
    \operatorname{Ext}^i_X(L_{X/Z},\mathcal M),
\]
and in particular
\[
    T^1(X/Z)=\operatorname{Ext}^1_X(L_{X/Z},\mathcal O_X)
\]
is the tangent space to the functor of first-order deformations of \(X\) over the
fixed target \(Z\); see \cite[\S 1.1--1.2]{Wehler1986}.

Apply the transitivity triangle for the composable morphisms
\[
    X \xrightarrow{\pi} Y \xrightarrow{i} Z.
\]
It gives a distinguished triangle
\[
    L\pi^*L_{Y/Z}
    \longrightarrow
    L_{X/Z}
    \longrightarrow
    L_{X/Y}
    \xrightarrow{+1}.
\]
Equivalently, in Wehler's notation, every morphism \(g:X\to Y\) over \(Z\)
induces a short exact sequence in the derived category, hence long exact
sequences in tangent cohomology; see \cite[\S 1.3(e)]{Wehler1986}. Applying
\(\mathbf R\mathcal Hom_X(-,\mathcal O_X)\) to the first arrow of the triangle gives
a canonical homomorphism
\[
    \operatorname{Ext}^1_X(L_{X/Z},\mathcal O_X)
    \longrightarrow
    \operatorname{Ext}^1_X(L\pi^*L_{Y/Z},\mathcal O_X).
\]
Thus we have a natural map
\[
    T^1(X/Z)
    \longrightarrow
    \operatorname{Ext}^1_X(L\pi^*L_{Y/Z},\mathcal O_X).
\]

\textcolor{blue}{Since \(i:Y\hookrightarrow Z\) is a closed immersion with ideal \(I\), the
relative cotangent complex is
\[
    L_{Y/Z}\simeq (I/I^2)[1];
\]
see, for instance, \cite[Ch.~III, Proposition~2.2.4]{Illusie1971} or
\cite[Proposition~3.2.9]{Sernesi2006}.}
We do not use the stronger assertion $$L_{Y/Z}\simeq (I/I^2)[1],$$ which would require the closed immersion $i:Y\hookrightarrow Z$ to be a local complete intersection immersion. Instead, we use only the first truncation of the cotangent complex. Locally on $Z$, the closed immersion $i:Y\hookrightarrow Z$ is given by a surjection of rings $A\to B=A/I$. For a surjective ring map $A\to B$ with kernel $I$, the cotangent complex satisfies $$H^0(L_{B/A})=0$$ and $$H^{-1}(L_{B/A})\simeq I/I^2;$$ see \cite[The Cotangent Complex, Lemma~92.11.2]{stacks-project}. Hence, after sheafification, one has $$H^0(L_{Y/Z})=0$$ and $$H^{-1}(L_{Y/Z})\simeq I/I^2.$$ By the canonical truncation triangle in the derived category $D(\mathcal O_Y)$, one has
$$
\tau_{\leq -2}L_{Y/Z}
\longrightarrow
L_{Y/Z}
\longrightarrow
\tau_{\geq -1}L_{Y/Z}
\xrightarrow{+1};
$$
see \cite[Derived Categories, Remark~13.12.4]{stacks-project}. Since $\tau_{\geq -1}L_{Y/Z}$ has only one nonzero cohomology sheaf, namely $I/I^2$ in degree $-1$, it follows that
$$
\tau_{\geq -1}L_{Y/Z}\simeq (I/I^2)[1].
$$
Set $$K:=\tau_{\leq -2}L_{Y/Z}.$$ Then $K\in D^{\leq -2}(\mathcal O_Y)$, while any $\mathcal O_Y$-module $\mathcal M$ is regarded as an object of $D^{\geq 0}(\mathcal O_Y)$. Therefore, by the standard Ext-vanishing for objects with bounded cohomological degrees, one has
$$
\operatorname{Ext}^n_Y(K,\mathcal M)=0
\qquad\text{for } n<2;
$$
see \cite[Derived Categories, Lemma~13.27.3]{stacks-project}. Applying the functor $\operatorname{Ext}^{\bullet}_Y(-,\mathcal M)$ to the truncation triangle gives an isomorphism
$$
\operatorname{Ext}^1_Y(L_{Y/Z},\mathcal M)
\simeq
\operatorname{Ext}^1_Y((I/I^2)[1],\mathcal M).
$$
Finally, using the shift convention for Ext groups, one obtains
$$
\operatorname{Ext}^1_Y((I/I^2)[1],\mathcal M)
=
\operatorname{Ext}^0_Y(I/I^2,\mathcal M)
=
\operatorname{Hom}_{\mathcal O_Y}(I/I^2,\mathcal M).
$$
Thus, for every $\mathcal O_Y$-module $\mathcal M$, and in particular for $\mathcal M=\pi_*\mathcal O_X$ or for a direct summand $\mathcal M=\mathcal E$, one has
$$
T^1(Y/Z,\mathcal M)
=
\operatorname{Ext}^1_Y(L_{Y/Z},\mathcal M)
\simeq
\operatorname{Hom}_{\mathcal O_Y}(I/I^2,\mathcal M).
$$

Therefore
\[
\begin{aligned}
    \operatorname{Ext}^1_X(L\pi^*L_{Y/Z},\mathcal O_X)
    &\simeq
    \operatorname{Ext}^1_X(L\pi^*((I/I^2)[1]),\mathcal O_X)  \\
    &\simeq
    \operatorname{Hom}_{D(X)}(L\pi^*(I/I^2),\mathcal O_X).
\end{aligned}
\]
By derived adjunction,
\[
    \operatorname{Hom}_{D(X)}(L\pi^*(I/I^2),\mathcal O_X)
    \simeq
    \operatorname{Hom}_{D(Y)}(I/I^2,R\pi_*\mathcal O_X).
\]
Since \(\pi\) is finite, \(R\pi_*\mathcal O_X=\pi_*\mathcal O_X\). Hence
\[
    \operatorname{Ext}^1_X(L\pi^*L_{Y/Z},\mathcal O_X)
    \simeq
    \operatorname{Hom}_{\mathcal O_Y}(I/I^2,\pi_*\mathcal O_X).
\]

Finally, using the given splitting
\[
    \pi_*\mathcal O_X\simeq \mathcal O_Y\oplus \mathcal E,
\]
we obtain
\[
    \operatorname{Hom}_{\mathcal O_Y}(I/I^2,\pi_*\mathcal O_X)
    \simeq
    \operatorname{Hom}_{\mathcal O_Y}(I/I^2,\mathcal O_Y)
    \oplus
    \operatorname{Hom}_{\mathcal O_Y}(I/I^2,\mathcal E).
\]
The composite is the required homomorphism \(\Theta=\Theta_1\oplus\Theta_2\).
\end{proof}

\begin{remark}
When \(X\) and \(Y\) are smooth and the deformation of \(\varphi=i\circ\pi\) is
locally trivial, \(T^1(X/Z)\) is represented by the usual global normal space
\(H^0(N_\varphi)\). In this case the map \(\Theta\) recovers the homomorphism
\[
    H^0(N_\varphi)
    \longrightarrow
    \operatorname{Hom}(\pi^*(I/I^2),\mathcal O_X)
    \simeq
    \operatorname{Hom}(I/I^2,\pi_*\mathcal O_X)
\]
used by Gonz\'alez. After the trace decomposition
\(\pi_*\mathcal O_X=\mathcal O_Y\oplus\mathcal E\), this gives the two components
\[
    H^0(N_\varphi)\to \operatorname{Hom}(I/I^2,\mathcal O_Y),
    \qquad
    H^0(N_\varphi)\to \operatorname{Hom}(I/I^2,\mathcal E),
\]
as in \cite[Proposition~3.7]{Gonzalez2005}; see also
\cite[Proposition~1.2]{GallegoGonzalezPurnaprajna2010}.
\end{remark}

\begin{theorem}[Central fiber of the image of a first-order deformation]
\label{thm:central-fiber-image-singular-source}
Let \(k\) be an algebraically closed field and let
\[
    X \xrightarrow{\pi} Y \xrightarrow{i} Z
\]
be morphisms of \(k\)-schemes, where \(\pi\) is finite and \(i\) is a closed immersion.
Let
\[
    \varphi=i\circ \pi:X\to Z,
\]
and let \(I=I_{Y/Z}\subset \mathcal O_Z\) be the ideal sheaf of \(Y\) in \(Z\).
Assume that
\[
    \pi_*\mathcal O_X\simeq \mathcal O_Y\oplus \mathcal E
\]
as an \(\mathcal O_Y\)-module.

Let
\[
    \Delta=\operatorname{Spec} k[\epsilon]/(\epsilon^2),
\]
and let
\[
    \widetilde\varphi:\widetilde X\longrightarrow Z\times \Delta
\]
be a first-order deformation of \(\varphi\) over the fixed target \(Z\), i.e. assume that
\(\widetilde X\) is flat over \(\Delta\), its central fiber is \(X\), and the restriction of
\(\widetilde\varphi\) to the central fiber is \(\varphi\).

Let
\[
    \widetilde Y:=\operatorname{Im}(\widetilde\varphi)\subset Z\times\Delta
\]
be the scheme-theoretic image of \(\widetilde\varphi\), and let
\[
    \widetilde Y_0:=\widetilde Y\times_{\Delta}\operatorname{Spec}k
    \subset Z
\]
be its central fiber.

Associated to \(\widetilde\varphi\) there is a natural homomorphism
\[
    \nu_{\widetilde\varphi}: I/I^2\longrightarrow \pi_*\mathcal O_X.
\]
Using the splitting
\[
    \pi_*\mathcal O_X\simeq \mathcal O_Y\oplus \mathcal E,
\]
write
\[
    \nu_{\widetilde\varphi}
    =
    \nu_1\oplus \nu_2,
\]
where
\[
    \nu_1:I/I^2\to \mathcal O_Y,
    \qquad
    \nu_2:I/I^2\to \mathcal E.
\]
Then the ideal sheaf of \(\widetilde Y_0\) in \(Z\) is
\[
    I_{\widetilde Y_0/Z}
    =
    \ker\left(
        I
        \longrightarrow
        I/I^2
        \xrightarrow{\nu_2}
        \mathcal E
    \right).
\]
In particular,
\[
    I^2\subset I_{\widetilde Y_0/Z}\subset I,
\]
and hence
\[
    Y\subset \widetilde Y_0\subset Y^{(1)}_Z,
\]
where \(Y^{(1)}_Z\) denotes the first infinitesimal neighborhood of \(Y\) in \(Z\).
\end{theorem}

\begin{proof}
The proof is local on \(Z\), so let
\[
    W=\operatorname{Spec}A\subset Z
\]
be an affine open subset and set
\[
    U=W\cap Y=\operatorname{Spec}(A/I),
\]
where, by abuse of notation, \(I\subset A\) denotes the ideal defining \(U\subset W\).
Let
\[
    V=\pi^{-1}(U)=\operatorname{Spec}B.
\]
Since \(\pi\) is finite, we have
\[
    B=\Gamma(U,\pi_*\mathcal O_X).
\]
Using the splitting \(\pi_*\mathcal O_X\simeq \mathcal O_Y\oplus \mathcal E\), after restricting to
\(U\) we may write
\[
    B\simeq A/I\oplus M,
    \qquad
    M=\Gamma(U,\mathcal E).
\]

Let
\[
    R=k[\epsilon]/(\epsilon^2).
\]
Then
\[
    W\times\Delta=\operatorname{Spec}(A\otimes_k R)
    =
    \operatorname{Spec}(A\oplus A\epsilon).
\]
The restriction of \(\widetilde X\) over \(V\) is affine, say
\[
    \widetilde V=\operatorname{Spec}\widetilde B.
\]
Since \(\widetilde X\) is flat over \(\Delta\) and has central fiber \(X\), the \(R\)-algebra
\(\widetilde B\) fits into a square-zero extension
\[
    0\longrightarrow B
    \xrightarrow{\iota}
    \widetilde B
    \xrightarrow{q}
    B
    \longrightarrow 0.
\]
Here \(q\) is reduction modulo \(\epsilon\), and \(\iota\) identifies \(B\) with
\(\epsilon\widetilde B\). Explicitly, if \(b\in B\) and \(\widetilde b\in \widetilde B\) is any lift of \(b\),
then
\[
    \iota(b)=\epsilon\widetilde b.
\]
This is independent of the chosen lift because \(\epsilon^2=0\). Thus
\[
    \ker(q)=\epsilon\widetilde B\simeq B,
    \qquad
    (\epsilon\widetilde B)^2=0.
\]

The morphism \(\widetilde\varphi\) gives a ring homomorphism
\[
    \widetilde\varphi^\#:
    A\oplus A\epsilon
    \longrightarrow
    \widetilde B.
\]
Modulo \(\epsilon\), this recovers the original morphism
\[
    \varphi^\#:A\to B.
\]
Since \(\varphi\) factors through \(Y\), the kernel of \(\varphi^\#:A\to B\) is \(I\), and the image
of \(\varphi^\#\) is the summand \(A/I\subset B=A/I\oplus M\).

We first construct the map
\[
    \nu_{\widetilde\varphi}: I/I^2\to B.
\]
If \(a\in I\), then
\[
    q\bigl(\widetilde\varphi^\#(a)\bigr)
    =
    \varphi^\#(a)
    =
    0.
\]
Therefore
\[
    \widetilde\varphi^\#(a)\in \ker(q)=\epsilon\widetilde B\simeq B.
\]
Define \(\lambda(a)\in B\) by the condition
\[
    \iota(\lambda(a))=\widetilde\varphi^\#(a).
\]
If \(a,b\in I\), then
\[
    \widetilde\varphi^\#(a),\widetilde\varphi^\#(b)\in \epsilon\widetilde B.
\]
Hence
\[
    \widetilde\varphi^\#(ab)
    =
    \widetilde\varphi^\#(a)\widetilde\varphi^\#(b)
    \in
    (\epsilon\widetilde B)^2
    =
    0.
\]
Thus \(\lambda(ab)=0\), so \(\lambda\) kills \(I^2\). Therefore \(\lambda\) descends to a homomorphism
\[
    \overline{\lambda}: I/I^2\to B.
\]
This is the local form of \(\nu_{\widetilde\varphi}\).

Now decompose
\[
    B=A/I\oplus M.
\]
Write
\[
    \overline{\lambda}=\nu_1\oplus \nu_2,
\]
where
\[
    \nu_1:I/I^2\to A/I,
    \qquad
    \nu_2:I/I^2\to M.
\]

Let \(\widetilde J\subset A\oplus A\epsilon\) be the ideal of the scheme-theoretic image
\(\widetilde Y\subset W\times\Delta\). Thus
\[
    \widetilde J=\ker(\widetilde\varphi^\#).
\]
We now compute \(\widetilde J\). Let \(a+a'\epsilon\in A\oplus A\epsilon\). If
\(a+a'\epsilon\in\widetilde J\), then reducing modulo \(\epsilon\) gives
\[
    \varphi^\#(a)=0,
\]
so \(a\in I\).

Assume now that \(a\in I\). Then
\[
    \widetilde\varphi^\#(a)=\iota(\lambda(a)).
\]
Also,
\[
    \widetilde\varphi^\#(a'\epsilon)
    =
    \epsilon\widetilde\varphi^\#(a')
    =
    \iota(\varphi^\#(a')).
\]
Therefore
\[
    \widetilde\varphi^\#(a+a'\epsilon)
    =
    \iota\bigl(\lambda(a)+\varphi^\#(a')\bigr).
\]
Since \(\iota\) is injective, this vanishes if and only if
\[
    \lambda(a)+\varphi^\#(a')=0
    \quad\text{in }B.
\]
Equivalently, writing \(B=A/I\oplus M\), this condition becomes
\[
    \nu_1(\overline a)+\overline{a'}=0
    \quad\text{in }A/I,
\]
and
\[
    \nu_2(\overline a)=0
    \quad\text{in }M,
\]
where \(\overline a\) denotes the class of \(a\) in \(I/I^2\), and
\(\overline{a'}\) denotes the class of \(a'\) in \(A/I\).

Thus
\[
    \widetilde J
    =
    \left\{
        a+a'\epsilon
        \ \middle|\
        a\in I,\ 
        \nu_1(\overline a)+\overline{a'}=0,\ 
        \nu_2(\overline a)=0
    \right\}.
\]

The ideal \(J\subset A\) of the central fiber \(\widetilde Y_0\subset W\) is
\[
    J=
    \frac{\widetilde J+(\epsilon)}{(\epsilon)}
    \subset A.
\]
Equivalently, \(a\in J\) if and only if there exists \(a'\in A\) such that
\[
    a+a'\epsilon\in \widetilde J.
\]
By the formula for \(\widetilde J\), this is equivalent to
\[
    a\in I,\qquad \nu_2(\overline a)=0,
\]
together with the existence of \(a'\in A\) satisfying
\[
    \overline{a'}=-\nu_1(\overline a).
\]
The last condition is automatic, because every element of \(A/I\) has a lift to \(A\).
Therefore
\[
    J
    =
    \{a\in I\mid \nu_2(\overline a)=0\}.
\]
Equivalently,
\[
    J=
    \ker\left(
        I\to I/I^2\xrightarrow{\nu_2}M
    \right).
\]
Sheafifying this local computation gives
\[
    I_{\widetilde Y_0/Z}
    =
    \ker\left(
        I_{Y/Z}
        \longrightarrow
        I_{Y/Z}/I_{Y/Z}^2
        \xrightarrow{\nu_2}
        \mathcal E
    \right).
\]

Since the map factors through \(I/I^2\), we have \(I^2\subset I_{\widetilde Y_0/Z}\).
Also \(I_{\widetilde Y_0/Z}\subset I\). Hence
\[
    I^2\subset I_{\widetilde Y_0/Z}\subset I,
\]
which is equivalent to
\[
    Y\subset \widetilde Y_0\subset Y^{(1)}_Z.
\]
This proves the theorem.
\end{proof}

\section{Deformations of double covers of cones over a rational normal curve}

\subsection{Canonical double covers of rational normal cones}

Let $e>4$, and let $Y=S(0,e)\subset \mathbb{P}^{e+1}$ be the rational normal cone of degree $e$. Equivalently, $Y\simeq \mathbb{P}(1,1,e)$. We denote by $p\in Y$ the vertex. The point $p$ is the unique singular point of $Y$, and analytically it is a cyclic quotient singularity of type $\frac{1}{e}(1,1)$.

Let $F$ denote the generator of the class group of $Y$, so that $\mathcal{O}_{Y}(F)=\mathcal{O}_{\mathbb{P}(1,1,e)}(1)$. The hyperplane class of the embedding $Y\subset \mathbb{P}^{e+1}$ is $H=eF$, or equivalently
$$
\mathcal{O}_{Y}(H)\simeq \mathcal{O}_{\mathbb{P}(1,1,e)}(e).
$$
The canonical sheaf of $Y$ is
$$
\omega_{Y}\simeq \mathcal{O}_{\mathbb{P}(1,1,e)}(-(e+2)).
$$

We consider a finite double cover $\pi:X\to Y$ whose trace decomposition is
$$
\pi_{*}\mathcal{O}_{X}=\mathcal{O}_{Y}\oplus \mathcal{E}.
$$
For the cover to be canonical with respect to the hyperplane class $H$, we require $\omega_{X}\simeq \pi^*\mathcal{O}_{Y}(H)$. The double-cover canonical formula gives
$$
\omega_{X}\simeq \pi^*(\omega_{Y}\otimes \mathcal{E}^{\vee}),
$$
where $\mathcal{E}^{\vee}:=\mathcal{H}om_{Y}(\mathcal{E},\mathcal{O}_{Y})$. Therefore the condition $\omega_{X}\simeq \pi^*\mathcal{O}_{Y}(H)$ forces $\omega_{Y}\otimes \mathcal{E}^{\vee}\simeq \mathcal{O}_{Y}(H)$, and hence
$$
\mathcal{E}\simeq \omega_{Y}(-H).
$$
Substituting the expressions above, we obtain
$$
\mathcal{E}
\simeq
\mathcal{O}_{\mathbb{P}(1,1,e)}(-(e+2)-e)
= \mathcal{O}_{\mathbb{P}(1,1,e)}(-(2e+2)).
$$

Thus the double cover is determined by the algebra $\mathcal{O}_{Y}\oplus\mathcal{E}$, together with a multiplication map
$$
\mathcal{E}\otimes_{\mathcal{O}_{Y}}\mathcal{E}\longrightarrow \mathcal{O}_{Y}.
$$
which is the same as a multiplication map 
$$
(\mathcal{E}\otimes_{\mathcal{O}_{Y}}\mathcal{E})^{\vee \vee}\longrightarrow \mathcal{O}_{Y}.
$$
Equivalently, this multiplication is determined by a section of the reflexive square $(\mathcal{E}^{\vee}\otimes_{\mathcal{O}_{Y}}\mathcal{E}^{\vee})^{\vee\vee}$. Since $\mathcal{E}^{\vee}\simeq \mathcal{O}_{\mathbb{P}(1,1,e)}(2e+2)$, we have
$$
(\mathcal{E}^{\vee}\otimes_{\mathcal{O}_{Y}}\mathcal{E}^{\vee})^{\vee\vee}
\simeq
\mathcal{O}_{\mathbb{P}(1,1,e)}(4e+4).
$$
Hence the branch data are given by a section
$$
s\in H^0\left(Y,\mathcal{O}_{\mathbb{P}(1,1,e)}(4e+4)\right).
$$
We denote the corresponding branch divisor by $B:=V(s)$. Thus
$$
B\in \left|\mathcal{O}_{\mathbb{P}(1,1,e)}(4e+4)\right|.
$$

Since $e>4$, every weighted homogeneous polynomial of degree $4e+4$ in coordinates of weights $1,1,e$ vanishes at the vertex $p=[0:0:1]$. Indeed, a monomial involving only the weight-$e$ coordinate would have degree divisible by $e$, while $4e+4\not\equiv 0\pmod e$ for $e>4$. Therefore every monomial of weighted degree $4e+4$ contains at least one of the weight-$1$ coordinates, and hence vanishes at $p$. Consequently, $p\in B$.

The trace-zero sheaf $\mathcal{E}=\mathcal{O}_{\mathbb{P}(1,1,e)}(-(2e+2))$ is rank-one reflexive. It is locally free at the vertex precisely when $e$ divides $2e+2$, equivalently when $e$ divides $2$. Since $e>4$, this does not occur. Thus $\mathcal{E}$ is not locally free at $p$, and the double cover is non-flat over the vertex in the sense of Alexeev--Pardini.

Away from $p$, the sheaf $\mathcal{E}$ is locally free, and the cover is locally given by an equation $z^2=f$. At the vertex, however, the cover must be understood through the reflexive algebra $\mathcal{O}_{Y}\oplus\mathcal{E}$ rather than through a line-bundle total space.

With this setup, the canonical divisor of $X$ satisfies $K_{X}=\pi^*H$. Since $H^2=e$ on the rational normal cone, one obtains
$$
K_{X}^2=2H^2=2e.
$$
Moreover,
$$
p_g(X)=h^0(X,\omega_{X})=h^0(Y,\mathcal{O}_{Y}(H))=e+2.
$$
Thus
$$
K_{X}^2=2e=2p_g(X)-4.
$$


\begin{proposition}[The singularity over the vertex]
\label{thm:singularity-canonical-double-cover-cone}
Let $e>4$, let $\pi:X\longrightarrow Y$
be the canonical double cover over $Y=S(0,e)\simeq \mathbb{P}(1,1,e)$ as in the set-up. Let $q\in X$ be the point lying over $p$. Then the following statements hold.

\begin{enumerate}
    \item On the local uniformizing cover
    $$
    \mathbb{A}^2_{u,v}
    \longrightarrow
    (Y,p)
    \simeq
    \left(
    \mathbb{A}^2_{u,v}/\mu_e,0
    \right),
    $$
    the pullback of the branch equation has the form
    $$
    f(u,v)
    =
    f_4(u,v)
    +
    f_{e+4}(u,v)
    +
    f_{2e+4}(u,v)
    +
    f_{3e+4}(u,v)
    +
    f_{4e+4}(u,v),
    $$
    where $f_d$ is homogeneous of ordinary degree $d$. For a general branch section, the binary quartic $f_4$ is squarefree, equivalently, its zero locus in $\mathbb{P}^1$ consists of four distinct points.

    \item The pullback of the germ $(X,q)$ to the local uniformizing cover is the hypersurface germ
    $$
    \widetilde X
    =
    \left\{
    w^2=f(u,v)
    \right\}
    \subset
    \mathbb{A}^3_{u,v,w},
    $$
    equipped with the $\mu_e$-action
    $$
    \zeta\cdot(u,v,w)
    =
    (\zeta u,\zeta v,\zeta^2w).
    $$
    Moreover,
    $$
    (X,q)
    \simeq
    (\widetilde X,0)/\mu_e.
    $$
   \item The minimal resolution
    $$
    \rho:\widehat X\longrightarrow X
    $$
    of the germ $(X,q)$ has a unique exceptional elliptic curve $C$ with
    $$
    C^2=-2e.
    $$
    Consequently, $(X,q)$ is a simple elliptic surface singularity. It is log canonical but not Kawamata log terminal, and the discrepancy of the exceptional elliptic curve is
    $$
    a(C,X)=-1.
    $$

    \item The surface $X$ is normal and Gorenstein. More precisely, its dualizing sheaf is invertible and satisfies
    $$
    \omega_X
    \simeq
    \pi^*\mathcal{O}_Y(H).
    $$
    In particular,
    $$
    K_X=\pi^*H
    $$
    is Cartier.
    \item For a general branch section, $X$ is smooth away from $q$. Consequently, $q$ is the unique singular point of $X$.
    \end{enumerate}
\end{proposition}


\begin{proof}[Proof of Part $(1)$]
Choose weighted homogeneous coordinates $[x:y:z]$ on
$Y=\mathbb P(1,1,e)$, with $\deg x=\deg y=1$ and $\deg z=e$. The vertex is
$p=[0:0:1]$. The affine chart $z\neq 0$ is analytically the quotient
$\mathbb A^2_{u,v}/\mu_e$, where $\mu_e$ acts by
$$
\zeta\cdot(u,v)=(\zeta u,\zeta v).
$$
Thus
$$
\mathbb A^2_{u,v}\longrightarrow (Y,p)
$$
is the local uniformizing cover.
Let $s(x,y,z)$ be the weighted homogeneous polynomial of degree $4e+4$
defining the branch section. Every monomial occurring in $s$ has the form
$x^ay^bz^j$, where $ a+b+ej=4e+4.$
Since $e>4$, one has $j\leq 4$.
Consequently, $s$ can be written uniquely as
$$
s(x,y,z)
\sum_{j=0}^{4}
z^j g_{4e+4-ej}(x,y),
$$
where each $g_{4e+4-ej}$ is an ordinary homogeneous polynomial in $x$ and
$y$ of degree $4e+4-ej$.

After pulling back to the uniformizing cover of the chart $z\neq 0$, the
coordinate $z$ becomes $1$, while $x$ and $y$ pull back to $u$ and $v$.
Therefore the pullback of the branch equation is
$$
f(u,v) =
f_4(u,v)
+
f_{e+4}(u,v)
+
f_{2e+4}(u,v)
+
f_{3e+4}(u,v)
+
f_{4e+4}(u,v),
$$
where $f_d$ is homogeneous of ordinary degree $d$. For a general branch divisor in $|\mathcal{O}_Y(4e+4)|$, $f_d$ is a general degree $d$ homogeneous polynomial in $u$ and $v$. Thus, for a
general branch section, $f_4$ is squarefree. Equivalently, its zero locus in
$\mathbb P^1$ consists of four distinct points.
\end{proof}
\begin{proof}[Proof of Part $(1)$]
For Part $2$, pull the double cover back along the local uniformizing morphism
$$
\mathbb A^2_{u,v}\longrightarrow (Y,p).
$$
Since the pullback of the branch section is the semi-invariant function $f(u,v)$ found in Part $1$, the pulled-back cover is the hypersurface germ
$$
\widetilde X =
\left\{
w^2=f(u,v)
\right\}
\subset
\mathbb A^3_{u,v,w}.
$$
It remains only to determine the induced $\mu_e$-action on $w$. The reflexive sheaf
$ \mathcal E \simeq \mathcal O_Y(-(2e+2))$
corresponds on the uniformizing cover to the module of semi-invariants of character
$-(2e+2)\equiv -2\pmod e.$
Equivalently, the dual generator defining the double-cover coordinate has character $2$. Thus the action on $w$ is $\zeta\cdot w=\zeta^2w.$ By construction, the quotient of this pulled-back double cover by $\mu_e$ recovers the original reflexive double-cover algebra near $p$. Hence

$$
(X,q)
\simeq
(\widetilde X,0)/\mu_e.
$$
This proves Part $2$.

\end{proof}

\begin{proof}[Proof of Part $(3)$]
Let $D\subset \mathbb A^2_{u,v}$ be the curve defined by $f(u,v)=0$, and let $o=(0,0)$. By Part $(1)$, one has $\operatorname{mult}_oD=4$, and the tangent cone of $D$ is the reduced divisor $f_4(u,v)=0$, consisting of four distinct lines.

Let $\sigma:S\to\mathbb A^2$ be the blowup of $o$, let $E\simeq\mathbb P^1$ be its exceptional curve, and let $D'$ be the strict transform of $D$. Then
$$
\sigma^*D=D'+4E
$$
and
$$
K_S=\sigma^*K_{\mathbb A^2}+E.
$$
Consequently,
$$
K_S+\frac12D'
= \sigma^*\left(K_{\mathbb A^2}+\frac12D\right)-E.
$$
Thus $E$ has discrepancy $-1$ with respect to the pair $\left(\mathbb A^2,\frac12D\right)$. This shows that $\widetilde{X}$ is log-canonical but not klt at $0$. \par
We now construct the corresponding resolution of $\widetilde X$. Let

$$
p:\widetilde X =\left\{w^2=f(u,v)\right\}\subset
\mathbb{A}^3_{u,v,w}\longrightarrow\mathbb A^2
$$
be the projection. \textcolor{blue}{Since the equation is monic in $w$}, the morphism $p$ is finite flat of degree $2$. Let $Z=p^{-1}(o)$ be the scheme-theoretic inverse image of the origin; it is defined on $\widetilde X$ by the ideal $(u,v)$ and is the length two subscheme $\operatorname{Spec}\displaystyle\frac{k[w]}{w^2}$ and is supported at $0 \in \widetilde X$.
Form the fiber product
$$
T:=\widetilde X\times_{\mathbb A^2}S.
$$
Since $p$ is flat, by the commuting of blow-ups and base change (see \cite{stacks-project}[Lemma $71.17.3$]), $T$ is the blowup of $\widetilde X$ along $Z$:

$$
T\simeq\operatorname{Bl}_Z(\widetilde X).
$$

Indeed, the lemma states that blowing up commutes with flat base change.

Let $\nu:\widehat{\widetilde X}\to T$ be the normalization, and denote by $\rho:\widehat{\widetilde X}\to\widetilde X$ the induced morphism. We claim that $\rho:\widehat{\widetilde X}\to\widetilde X$ is a resolution. For this it is enough to verify that $\widehat{\widetilde X}$ is smooth. Since the map $\rho:\widehat{\widetilde X}\to\widetilde X$ is birational, it is enough to verify smooth above the point $0 \in \widetilde X$. But this locus is the same as the scheme theoretic inverse image of the exceptional locus of $S \to \mathbb{A}^2$ along the map $\widehat{\widetilde X} \to S$. We verify that $\widehat{\widetilde X}$ is smooth above the exceptional locus. On the chart of $S$ given by $v=ux$, write
$$
f(u,ux)= u^4\left(f_4(1,x)+u^ef_{e+4}(1,x)+u^{2e}f_{2e+4}(1,x)+u^{3e}f_{3e+4}(1,x)+u^{4e}f_{4e+4}(1,x)\right).
$$
Let $$F(u,x) =  f_4(1,x)+u^ef_{e+4}(1,x)+u^{2e}f_{2e+4}(1,x)+u^{3e}f_{3e+4}(1,x)+u^{4e}f_{4e+4}(1,x)$$ denote the strict transform $D'$ of $D$ under the blow-up. One has $F(0,x)=f_4(1,x).$
On this chart, the fiber product $T$ is given by

$$
\{w^2=u^4F(u,x)\} \subset \mathbb{A}^3_{u,x,w}.
$$
In its function field, the element $\eta=w/u^2$ is integral, since it satisfies $\eta^2=F(u,x)$. Hence the normalization is given locally by
$$
\{\eta^2=F(u,x)\} \subset \mathbb{A}^3_{u,x,\eta}
$$
Along the exceptional divisor $u=0$, this equation becomes
$$
\eta^2=f_4(1,x).
$$
If this hypersurface is singular at a point $(0,x_0, \eta_0, )$, all partials have to vanish which implies $\eta_0 = f(1,x_0) = \frac{d}{dx}f(1,x)|_{x=x_0} = 0$. The analogous calculation on the other blowup chart gives the same conclusion. But this means that the strict transform of $f$ meets the exceptional divisor at a point with multiplicity at least two, which contradicts that $f(u,v)$ is a general quartic.
Therefore $\widehat{\widetilde X}$ is smooth above the exceptional locus, and, after restricting to the germ at the origin, $\rho$ is a resolution.

Let $h:\widehat{\widetilde X}\to S$ be the induced double cover, and let $C=h^{-1}(E)$. The local equation above shows that $E$ (given by $(u = 0)$) is not a component of the branch divisor of $h$; the branch divisor is $D'$ whose local equation is given by $F(u,x)$. Moreover, the restriction
$$
C\longrightarrow E\simeq\mathbb P^1
$$
is the double cover branched along the four distinct points $D'\cap E$ (this double cover is locally defined by $\eta^2=f_4(1,x)$). Hence $C$ is a smooth elliptic curve.
Since $E$ is not a branch component, one has $h^*E=C$. By the projection formula,
$C^2=(h^*E)^2=2E^2=-2.$ Thus the exceptional locus of $\rho$ is the single smooth elliptic curve $C$ with $C^2=-2$.
Finally, the double-cover canonical formulas give

$$
K_{\widetilde X} =
p^*\left(K_{\mathbb A^2}+\frac12D\right)
$$

and

$$
K_{\widehat{\widetilde X}} =
h^*\left(K_S+\frac12D'\right).
$$

Using the discrepancy calculation for the pair and the equality $h^*E=C$, we obtain

$$
K_{\widehat{\widetilde X}}=
\rho^*K_{\widetilde X}-C.
$$

Therefore the discrepancy of the exceptional elliptic curve over $\widetilde X$ is $-1$.
\end{proof}
\begin{proposition}
Let $e>4$, let $\pi:X\longrightarrow Y$
be the canonical double cover over $Y=S(0,e)\simeq \mathbb{P}(1,1,e)$ as in the set-up. Let $q\in X$ be the point lying over $p$.
\begin{enumerate}
    \item Let
    $$
    C
    =
    \left\{
    w^2=f_4(u,v)
    \right\}
    \subset
    \mathbb{P}(1,1,2),
    \qquad
    L:=\mathcal{O}_C(1).
    $$
    Then $C$ is a smooth elliptic curve and
    $$
    \deg L=2.
    $$
    The germ $(X,q)$ is analytically isomorphic to the vertex of the affine cone
    $$
    \operatorname{Spec}
    \left(
    \bigoplus_{m\geq 0}
    H^0(C,L^{em})
    \right).
    $$
    Equivalently, $(X,q)$ is the affine cone over the embedding of $C$ defined by the complete linear system
    $$
    |L^e|.
    $$
    Since
    $$
    \deg L^e=2e,
    $$
    this realizes $C$ as an elliptic normal curve of degree $2e$ in
    $$
    \mathbb{P}^{2e-1}.
    $$

    
\end{enumerate}
\end{proposition}




\begin{proposition}\label{prop:hom-conormal-trace-zero-cone}
Let $Y=S(0,e)\subset\mathbb{P}^{e+1}$ be the rational normal cone of degree $e$, and let $H=\mathcal{O}_Y(1)$ denote the hyperplane class of this embedding. Equivalently, identify
$$
Y\simeq \mathbb{P}(1,1,e),
$$
so that
$$
H\simeq \mathcal{O}_{\mathbb{P}(1,1,e)}(e).
$$
Let $\mathcal{I}\subset\mathcal{O}_{\mathbb{P}^{e+1}}$ be the ideal sheaf of $Y$. Set
$$
\mathcal{E}:=\omega_Y(-H).
$$
Then
$$
\operatorname{Hom}_Y(\mathcal{I}/\mathcal{I}^2,\mathcal{E})=0.
$$
\end{proposition}

\begin{proof}
The homogeneous ideal of the rational normal cone $Y=S(0,e)\subset\mathbb{P}^{e+1}$ is generated by quadrics. Therefore there is a surjection
$$
\mathcal{O}_Y(-2H)^{\oplus N}\twoheadrightarrow \mathcal{I}/\mathcal{I}^2
$$
for some integer $N$. Applying $\operatorname{Hom}_Y(-,\mathcal{E})$ gives an injection
$$
\operatorname{Hom}_Y(\mathcal{I}/\mathcal{I}^2,\mathcal{E})
\hookrightarrow
\operatorname{Hom}_Y(\mathcal{O}_Y(-2H)^{\oplus N},\mathcal{E}).
$$
The right-hand side is
$$
H^0(Y,\mathcal{E}(2H))^{\oplus N}.
$$
Thus it is enough to show that
$$
H^0(Y,\mathcal{E}(2H))=0.
$$

Now
$$
\mathcal{E}(2H)=\omega_Y(-H+2H)=\omega_Y(H).
$$
Under the identification $Y\simeq\mathbb{P}(1,1,e)$, one has
$$
\omega_Y\simeq\mathcal{O}_{\mathbb{P}(1,1,e)}(-(e+2)).
$$
Since $H\simeq\mathcal{O}_{\mathbb{P}(1,1,e)}(e)$, it follows that
$$
\mathcal{E}(2H)
\simeq
\mathcal{O}_{\mathbb{P}(1,1,e)}(-(e+2)+e)
= \mathcal{O}_{\mathbb{P}(1,1,e)}(-2).
$$
But
$$
H^0(\mathbb{P}(1,1,e),\mathcal{O}_{\mathbb{P}(1,1,e)}(-2))=0.
$$
Therefore
$$
H^0(Y,\mathcal{E}(2H))=0.
$$
Hence
$$
\operatorname{Hom}_Y(\mathcal{I}/\mathcal{I}^2,\mathcal{E})=0.
$$
\end{proof}



\bibliographystyle{plain}
\begin{thebibliography}{KMM92}
\color{black}



\bibitem{AHL10}
Artebani, Michela; Hausen, Jürgen; Laface, Antonio
\emph{On Cox rings of K3 surfaces}
Compos. Math. 146 (2010), no. 4, 964–998.

\bibitem{ABS17}
Arbarello, Enrico; Bruno, Andrea; Sernesi, Edoardo
\emph{On hyperplane sections of K3 surfaces.}
Algebr. Geom. 4 (2017), no. 5, 562–596.


\bibitem{AltmannStevens}
K.~Altmann and J.~Stevens,
\newblock \emph{Cotangent cohomology of rational surface singularities},
\newblock Invent. Math. \textbf{138} (1999), no.~1, 163--181.

\bibitem{BE95}
Bayer, Dave; Eisenbud, David
\emph{Ribbons and their canonical embeddings}
Trans. Amer. Math. Soc. 347 (1995), no. 3, 719–756.


\bibitem{Bea02}
Beauville, A. \emph{Fano threefolds and $K3$ surfaces}
https://arxiv.org/abs/math/0211313
%\bibitem[Bel24]{Bel24}
%Belmans, Peter
%\emph{Fanography: a tool to visually study the geography of Fano 3-folds.}
%https://www.fanography.info

\bibitem{BGM24} Bangere, Purnaprajna; Gallego, Francisco Javier; Mukherjee; Jayan \emph{Projective smoothing of varieties with simple normal crossings.}(https://arxiv.org/abs/2403.04167)

\bibitem{BKM25}
Bayer, Arend;  Kuznetsov, Alexander;  Macr\'i, Emanuele 
\emph{Mukai models of Fano varieties}
arXiv:2501.16157

\bibitem{BM87}
Beauville, A.; Mérindol, J.-Y.
\emph{Sections hyperplanes des surfaces $K3$.[Hyperplane sections of $K3$ surfaces]}
Duke Math. J. 55 (1987), no. 4, 873–878.

\bibitem{BMR21}
Bangere, Purnaprajna; Mukherjee, Jayan; Raychaudhury, Debaditya
\emph{K3 carpets on minimal rational surfaces and their smoothings.}
Internat. J. Math.32(2021), no.6, Paper No. 2150032, 20 pp.

\bibitem{BM24} 
Bangere, Purnaprajna; Mukherjee, Jayan  
\emph{Extendability of projective varieties via degeneration to ribbons with applications to Calabi-Yau threefolds} (arXiv:2409.03960)

\bibitem{BM25}
Bangere, Purnaprajna; Mukherjee, Jayan  
\emph{Extendability of general prime $K3$ surfaces via degeneration to ribbons} (in progress)



%\bibitem[BGM24]{BGM24} Bangere, Purnaprajna; Gallego, Francisco Javier; Mukherjee; Jayan \emph{Projective degenerations of $K3$ surfaces to a union of Hirzebruch surfaces not embedded as scrolls.} (in preparation)

%\bibitem[BGM24b]{BGM24b} Bangere, Purnaprajna; Gallego, Francisco Javier; Mukherjee; Jayan \emph{Projective smoothing of simple normal crossing projective bundles.} (in preparation)

%\bibitem[BGM24c]{BGM24c} Bangere, Purnaprajna; Gallego, Francisco Javier; Mukherjee; Jayan \emph{Projective smoothing of simple normal crossing complete intersections.} (in preparation)

\bibitem{CD22}
Ciliberto, Ciro; Dedieu, Thomas
\emph{$K3$ curves with index $k>1$.}
Boll. Unione Mat. Ital. 15 (2022), no. 1-2, 87–115.

\bibitem{CD24}
Ciliberto, Ciro; Dedieu, Thomas
\emph{Double covers and extensions.}
Kyoto J. Math. 64 (2024), no. 1, 75–94.


\bibitem{CDGK20}
Ciliberto, C; Dedieu, T;  Galati, C.; Knutsen, A. L.  
\emph{Moduli of curves on Enriques surfaces.}
Adv. Math. 365 (2020), 107010, 42 pp. 6, 13, 14

\bibitem{CDS20}
Ciliberto, Ciro; Dedieu, Thomas; Sernesi, Edoardo
\emph{Wahl maps and extensions of canonical curves and $K3$ surfaces.}
J. Reine Angew. Math. 761 (2020), 219–245.

\bibitem{CLM93}
Ciliberto, C; Lopez, A; Miranda, R
\emph{Projective degenerations of $K3$ surfaces, Gaussian maps and Fano threefolds}
Invent. Math. 114, 641 667 (1993)

\bibitem{CLM94}
Ciliberto, Ciro; Lopez, Angelo Felice; Miranda, Rick
\emph{On the corank of Gaussian maps for general embedded $K3$ surfaces}
arXiv:alg-geom/9411008

\bibitem{CLM98}
Ciliberto, Ciro; Lopez, Angelo Felice; Miranda, Rick
\emph{Classification of varieties with canonical curve section via Gaussian maps on canonical curves.}
Amer. J. Math.   120 (1998), no. 1, 1–21.

\bibitem{CLM98c}
Ciliberto, Ciro; Lopez, Angelo Felice; Miranda, Rick
\emph{Corrigendum to ``Classification of varieties with canonical curve section via Gaussian maps on canonical curves''.}
Amer. J. Math. 143 (2021), no. 6, 1661–1663.

\bibitem{CilibertoLopezMiranda}
C.~Ciliberto, A.~F.~Lopez, and R.~Miranda,
\newblock \emph{Some remarks on the obstructedness of cones over curves of low genus},
\newblock in \emph{Higher Dimensional Complex Varieties (Trento, 1994)},
de Gruyter, Berlin, 1996, pp.~167--182.

\bibitem{CM90}
Ciliberto, Ciro; Miranda, Rick
\emph{On the Gaussian map for canonical curves of low genus.}
Duke Math. J. 61 (1990), no. 2, 417–443.

\bibitem{DM92}
Duflot, Jeanne; Miranda, Rick
\emph{The Gaussian map for rational ruled surfaces.}
Trans. Amer. Math. Soc. 330 (1992), no. 1, 447–459.

\bibitem{DM24}
Deopurkar, Anand; Mukherjee, Jayan
\emph{Syzygies of canonical ribbons on higher genus curves}
arXiv:2412.05500

%\color{black}
%\bibitem[Deb01]{Debarre} Debarre, O., \emph{Higher-dimensional algebraic geometry}, 
%Universitext Springer--Verlag, New York, 2001.
%\color{black}

\bibitem{DP12}
Duflot, Jeanne; Peters, Pamela L.
\emph{Gaussian maps for double covers of toric surfaces.}
Rocky Mountain J. Math. 42 (2012), no. 5, 1471–1520.












%\bibitem[Fri83]{F83}
%Friedman, Robert; 
%\emph{Global smoothings of varieties with normal crossings.} 
%Ann. of Math. (2) 118 (1983), no. 1, 75--114.

 

\color{black}

%\bibitem[Fuj14a]{KFujita} Fujita, Kento, \emph{Simple normal crossing Fano varieties and log Fano manifolds.}
%Nagoya Math. J. 214 (2014), 95-123.

%\bibitem[Fuj14b]{Fujino} Fujino, Osamu.
%\emph{Fundamental theorems for semi log canonical pairs.} Algebr. Geom. 1 (2014), no.2, 194-228.

%\color{black}

\bibitem{GGP08}
Gallego, Francisco Javier; González, Miguel; Purnaprajna, Bangere P., 
\emph{$K3$ double structures on Enriques surfaces and their smoothings.}
J. Pure Appl. Algebra 212 (2008), no. 5, 981--993

\bibitem{GGP13}
Gallego, F.J., González, M., Purnaprajna, B.P., 
\emph{An infinitesimal condition to smooth ropes.}
Rev Mat Complut 26, 253–269 (2013).

\bibitem{Gon06}
  Gonz\'alez, Miguel.
  \emph{Smoothing of ribbons over curves}.
  J. Reine Angew. Math. 591 (2006), 201--235.

\bibitem{GP97}
Gallego, F.J., Purnaprajna, B.P.
\emph{Degenerations of $K3$ surfaces in projective space}
Transactions of the American Mathematical Society
Volume 349, Number 6, June 1997, Pages 2477–2492



\color{black}

  %\bibitem[Gro61]{EGA} A. Grothendieck, \emph{\'El\'ements 
%de g\'eom\'etrie 
%alg\'ebrique. III. \'Etude cohomologique des faisceaux 
%coh\'erents. I.}, Inst. 
%Hautes \'Etudes Sci. Publ. Math.  
%{11}, (1961).

\bibitem{Hartshorne}  
Hartshorne, R; \emph{Algebraic geometry}, 
Grad. Texts in Math., No. 52
Springer-Verlag, New York-Heidelberg, 1977

\color{black}

\bibitem{Hor74}
Horikawa, Eiji
\emph{On deformations of holomorphic maps. II.}
J. Math. Soc. Japan 26 (1974), 647–667.



\bibitem{I77}
Iskovskih, V. A.
\emph{Fano threefolds. I.}
(Russian)Izv. Akad. Nauk SSSR Ser. Mat.41 (1977), no.3, 516–562, 717.

\bibitem{I78}
Iskovskih, V. A.
\emph{Fano threefolds. II.}
(Russian)Izv. Akad. Nauk SSSR Ser. Mat.42 (1978), no.3, 506–549.



%\bibitem[Kim17]{Ki17}
%Kimura, Y.;
%\emph{Structure of stable degeneration of $K3$ surfaces into pairs of rational elliptic surfaces}
%JHEP03(2018)045, %https://doi.org/10.48550/arXiv.1710.04984
\bibitem{K01}
Knutsen, Andreas Leopold
\emph{On $k$th-order embeddings of $K3$ surfaces and Enriques surfaces.}
Manuscripta Math. 104 (2001), no. 2, 211–237.

\bibitem{K20}
Knutsen, Andreas Leopold
\emph{Global sections of twisted normal bundles of $K3$ surfaces and their hyperplane sections.}
Atti Accad. Naz. Lincei Rend. Lincei Mat. Appl. 31 (2020), no. 1, 57–79.


\bibitem{KLM11}
Knutsen, Andreas Leopold; Lopez, Angelo Felice; Muñoz, Roberto
\emph{On the extendability of projective surfaces and a genus bound for Enriques-Fano threefolds.}
J. Differential Geom.88(2011), no.3, 485–518.

\bibitem{KL81}
Kleppe, J.; 
\emph{The Hilbert Flag Scheme, its properties, and its connection with the Hilbert Scheme. Applications to curves in three-space. Thesis. University of Oslo, Preprint No. 5 (1981)}
\bibitem{KL87}
Kleppe, J.;
\emph{Non-reduced components of the Hilbert scheme of smooth space curves. In: Space Curve.}
Rocca di papa, 1985. Ghione, F., Peskine, C., Sernesi, E. (eds.) Springer Lect. Notes Math. No. 1266 (1987), 181-207

\bibitem{HR} M.~Hochster and J.\,L.~Roberts, \emph{Rings of invariants of reductive groups acting on regular rings are Cohen--Macaulay}, Advances in Math.\ \textbf{13} (1974), 115--175.
 
\bibitem{Illusie} L.~Illusie, \emph{Complexe cotangent et d\'eformations I, II}, Lecture Notes in Math.\ \textbf{239}, \textbf{283}, Springer-Verlag, 1971--1972.
 
\bibitem{Saito} M.~Saito, \emph{Mixed Hodge complexes on algebraic varieties}, Math.\ Ann.\ \textbf{316} (2000), 283--331.

\bibitem{stacks-project}
The Stacks Project Authors,
\emph{The Stacks Project},
\url{https://stacks.math.columbia.edu}
 
\bibitem{Steenbrink} J.\,H.\,M.~Steenbrink, \emph{Mixed Hodge structure on the vanishing cohomology}, in: Real and Complex Singularities, Sijthoff and Noordhoff, 1977, 525--563.
 
\bibitem{Wehler} J.~Wehler, \emph{Cyclic coverings: deformation and Torelli theorem}, Math.\ Ann.\ \textbf{274} (1986), 443--472.

  




%\bibitem[Kon85]{Ko85}
%Kondo, S
%\emph{Type II degenerations of $K3$ surfaces}
%Nagoya Math. J. Vol. 99 (1985), 11-30



%\bibitem[Kul77]{Ku77}
%Kulikov, V. 
%\emph{Degenerations of K3 surfaces and Enriques surfaces} Math. USSR Izvestiya 11 (1977), 957-989.





\color{black}



\bibitem{Lop23}
Lopez, Angelo Felice
\emph{On the extendability of projective varieties: a survey.The art of doing algebraic geometry, 261–291. With an appendix by Thomas Dedieu}
Trends Math. Birkhäuser/Springer, Cham, [2023], ©2023 ISBN:978-3-031-11937-8
ISBN:978-3-031-11938-5


%\bibitem[LR12]{LR}
%Liu, Wenfei; Rollenske, Sönke. \emph{Two-dimensional semi-log-canonical hypersurfaces}, 
%Matematiche (Catania) 67 (2012), no. 2, 185–202.

\bibitem{Lvo92} 
L\'vovsky, S.
\emph{Extensions of projective varieties and deformations. I, II.}
Michigan Math. J.39(1992), no.1, 41–51, 65--70.





\bibitem{MM82}
Mori, Shigefumi; Mukai, Shigeru
\emph{Classification of Fano 3-folds with $B2 \geq 2$.}
Manuscripta Math.36(1981/82/1982), no.2, 147–162.


\bibitem{M88}
Mukai, S
\emph{Curves, K3 Surfaces and Fano 3-folds of Genus $\leq$ 10}
Algebraic Geometry and Commutative Algebra, Academic Press,
1988, Pages 357-377.

  \color{black}
  


  \color{black}
%\bibitem[Par91]{Par91}
%Pardini, Rita
%\emph{Abelian covers of algebraic varieties.}
%J. Reine Angew. Math.417(1991), 191–213.




%\bibitem[PP81]{PP81}
%Persson, U; Pinkham, H 
%\emph{Degenerations of surfaces with trivial canonical bundle}
%Ann. of Math. 113 (1981), 45-66

\bibitem{Pinkham}
H.~C.~Pinkham,
\newblock \emph{Deformations of algebraic varieties with \(G_m\) action},
\newblock Ast\'erisque \textbf{20}, Soci\'et\'e Math\'ematique de France, Paris, 1974.


\bibitem{Ser}
Sernesi, Edoardo.
\emph{Deformations of Algebraic Schemes.}
Grundlehren der Mathematischen Wissenschaften [Fundamental Principles of Mathematical Sciences], 334. Springer-Verlag, Berlin, 2006.

\bibitem[Stacks]{stacks-project}
The Stacks Project Authors,
\emph{Stacks Project}.
\url{https://stacks.math.columbia.edu}.

\bibitem{StacksProject}
The Stacks Project Authors,
\emph{Stacks Project},
\url{https://stacks.math.columbia.edu}.

\bibitem{Tot20}
Totaro, Burt
\emph{Bott vanishing for algebraic surfaces.}
Trans. Amer. Math. Soc. 373 (2020), no. 5, 3609–3626.
%\bibitem[Sei92]{Sei92}
%Seiler, Wolfgang K.
%\emph{Deformations of ruled surfaces.}
%J. Reine Angew. Math.426(1992), 203–219.

%\bibitem[SP]{SP}
%Stacks project authors
%\emph{The Stacks project}
%\url{https://stacks.math.columbia.edu}, 2023

%\bibitem[Tho13]{T13}
%Thompson, A;
%\emph{Degenerations of $K3$ surfaces of degree two}
%Transactions of the American Mathematical Society
%Volume 366, Number 1, January 2014, Pages 219–243
%S 0002-9947(2013)05759-5


\bibitem{V92}
Voisin, Claire
\emph{Sur l'application de Wahl des courbes satisfaisant la condition de Brill-Noether-Petri.[On the Wahl mapping of curves satisfying the Brill-Noether-Petri condition]}
Acta Math. 168 (1992), no. 3-4, 249–272.






\bibitem{W87}
Wahl, J. M. 
\emph{The Jacobian algebra of a graded Gorenstein singularity.}
Duke Math. J. 55 (1987), no. 4,
843-871

\bibitem{Zak91}
Zak, F. L.
\emph{Some properties of dual varieties and their applications in projective geometry.}
Algebraic geometry (Chicago, IL, 1989), 273–280.
Lecture Notes in Math., 1479
Springer-Verlag, Berlin, 1991

\bibitem{Wehler1986}
J.~Wehler,
\newblock Cyclic coverings: deformation and Torelli theorem,
\newblock \emph{Math. Ann.} \textbf{274} (1986), 443--472.

\bibitem{Gonzalez2005}
M.~Gonz\'alez,
\newblock Smoothing of ribbons over curves,
\newblock \emph{J. Reine Angew. Math.} \textbf{591} (2006), 201--235.

\bibitem{GallegoGonzalezPurnaprajna2010}
F.~J. Gallego, M.~Gonz\'alez, and B.~P. Purnaprajna,
\newblock Deformation of canonical morphisms and the moduli of surfaces of general type,
\newblock \emph{Invent. Math.} \textbf{182} (2010), 1--46.

\bibitem{Illusie1971}
L.~Illusie,
\newblock \emph{Complexe cotangent et d\'eformations I},
\newblock Lecture Notes in Mathematics, Vol.~239,
\newblock Springer-Verlag, Berlin--New York, 1971.

\bibitem{Sernesi2006}
E.~Sernesi,
\newblock \emph{Deformations of Algebraic Schemes},
\newblock Grundlehren der mathematischen Wissenschaften, Vol.~334,
\newblock Springer-Verlag, Berlin, 2006.




 
\end{thebibliography}

\end{document}






\begin{remark}[2.4]
The content of Lemma~2.3 is a purely cohomological fact of sheaves on $X$.
The proof given in \cite[4.2]{Hor74}, in the context of complex manifolds, is still valid
in the case of a homomorphism of quasi--coherent sheaves
\[
\mathcal A \xrightarrow{\,F\,} \mathcal B
\]
with kernel and cokernel, respectively, $\mathcal K$ and $\mathcal N$, on a noetherian separated scheme $X$.
Let $\mathcal V$ be an open cover of $X$. We set
\begin{equation}\tag{2.4.1}
D(X,F;\mathcal V)
=
\frac{\Bigl\{(g,\rho)\in C^{0}(\mathcal V,\mathcal B)\times Z^{1}(\mathcal V,\mathcal A)\ \Big|\ 
\delta g = F\rho\Bigr\}}
{\Bigl\{(Fh,\delta h)\ \Big|\ h\in C^{0}(\mathcal V,\mathcal A)\Bigr\}}
\end{equation}
and
\begin{equation}\tag{2.4.2}
D(X,F)=\varinjlim_{\mathcal V}\, D(X,F;\mathcal V),
\end{equation}
where the direct limit is taken under refinement of coverings. Then for every open affine cover
$\mathcal V$ the natural map $D(X,F;\mathcal V)\to D(X,F)$ is an isomorphism.
Two exact sequences, like in Lemma~2.3, are then obtained:
\[
H^{0}(\mathcal A)\longrightarrow H^{0}(\mathcal B)\longrightarrow D(X,F)
\longrightarrow H^{1}(\mathcal A)\longrightarrow H^{1}(\mathcal B),
\]
\[
0\longrightarrow H^{1}(\mathcal K)\longrightarrow D(X,F)\longrightarrow H^{0}(\mathcal N)
\longrightarrow H^{2}(\mathcal K).
\]
\end{remark}